\begin{pastexam}[label={px:nccu-107-dsa-07}]{政大資科所 107 學年度 第 7 題}
（問答）

(15\%) A contiguous subsequence of a list $S$ is a subsequence made up of consecutive elements of
$S$. For instance, if $S$ is 5, 15, $-30$, 10, $-5$, 40, 10, then 15, $-30$, 10 is a contiguous
subsequence but 5, 15, 40 is not. Design a polynomial-time algorithm for the following task:\\
Input: A list $S$ of numbers $a_1,a_2,\ldots,a_n$.\\
Output: The contiguous subsequence of $S$ of maximum sum.\\
Note that a subsequence of length zero has sum zero.\\
For the preceding example, the answer would be 10, $-5$, 40, 10, with a sum of 55.
\end{pastexam}

\begin{pastsoln}
\textbf{解讀}。 題目允許長度 $0$ 的子序列，和是 $0$。 所以全部都是負數時，答案是空的子序列（和 $0$），
這和資結書第~17.2~節「至少一個數」的約定不同。 輸出用位置 $i..j$ 表示 $a_i, \ldots, a_j$，空的
子序列寫成 $j = i - 1$。 多段一樣大時，回傳最先達到最大和的那一段。

\textbf{想法}。 改寫資結書第~17.2.3~節的線性做法。 令
\[
F[j] = \text{結尾在 } j \text{ 的子序列（可以是空的）的最大和},
\]
$F[0] = 0$。 結尾在 $j$ 的非空子序列，是「結尾在 $j - 1$ 的子序列（可以是空的）」再接上 $a_j$，最大
是 $F[j-1] + a_j$；再和空的比，$F[j] = \max(0,\ F[j-1] + a_j)$。 答案是所有「結尾在 $j$ 的非空最佳
和」$F[j-1] + a_j$ 和空序列的 $0$ 之中最大的。

\begin{pseudo}{最大和的連續子序列（允許空的）}
\Function{Max-Contiguous}{$a, n$}
    \State $\mathit{best} \gets 0$；$\mathit{bi} \gets 1$；$\mathit{bj} \gets 0$ \Comment{先記空的子序列}
    \State $\mathit{cur} \gets 0$；$\mathit{start} \gets 1$
    \For{$j \gets 1$ \To $n$}
        \State $\mathit{cur} \gets \mathit{cur} + a_j$
        \If{$\mathit{cur} > \mathit{best}$}
            \State $\mathit{best} \gets \mathit{cur}$；$\mathit{bi} \gets \mathit{start}$；$\mathit{bj} \gets j$
        \EndIf
        \If{$\mathit{cur} < 0$}
            \State $\mathit{cur} \gets 0$；$\mathit{start} \gets j + 1$ \Comment{接上前面只會更差}
        \EndIf
    \EndFor
    \State \Return $(\mathit{bi}, \mathit{bj}, \mathit{best})$
\EndFunction
\end{pseudo}

\textbf{正確性}。 不變量：第 4 行每一輪開始時，(i) $\mathit{cur} = F[j-1] = a_{\mathit{start}} + \cdots
+ a_{j-1}$（$\mathit{start} = j$ 時是空的、和為 $0$）；(ii) $(\mathit{bi}, \mathit{bj}, \mathit{best})$
是 $a_1, \ldots, a_{j-1}$ 裡（可以是空的）和最大的連續子序列。
\begin{itemize}
\item \textbf{初始}：$j = 1$，$\mathit{cur} = 0 = F[0]$，只有空的子序列，(i)(ii) 都成立。
\item \textbf{維持}：第 5 行之後 $\mathit{cur} = F[j-1] + a_j = a_{\mathit{start}} + \cdots + a_j$，是
結尾在 $j$ 的非空子序列的最大和。 $a_1..a_j$ 裡的子序列，不是在 $a_1..a_{j-1}$ 裡，就是結尾在 $j$
的非空子序列，所以第 6--7 行之後 (ii) 對 $j$ 成立。 第 8--9 行把 $\mathit{cur}$ 變成
$\max(0, \mathit{cur}) = F[j]$：小於 $0$ 時換成空的子序列，從 $j + 1$ 重新開始。 (i) 對下一輪成立。
\item \textbf{結束}：$j = n + 1$，由 (ii)，回傳的是整個 $S$ 裡和最大的連續子序列。
\end{itemize}

\textbf{時間}。 迴圈 $n$ 輪，每輪常數個動作，$\Theta(n)$；額外空間 $\Theta(1)$。 題目只要求多項式時間，
兩兩枚舉所有 $i \le j$ 邊走邊加的 $\Theta(n^2)$ 或資結書第~17.2~節的分治 $\Theta(n \lg n)$ 也都符合，
但線性做法最好，而且不可能更快：每個數都可能是答案的一部分，至少要讀過一次。

\textbf{追蹤題目的例子}。 $S = 5, 15, -30, 10, -5, 40, 10$。 表中「第 5 行後」是加上 $a_j$ 之後的
$\mathit{cur}$，「本輪結束」是第 8--9 行之後的值。
\begin{center}
\begin{tabular}{ccccl}
\toprule
$j$ & $a_j$ & 第 5 行後 & 本輪結束 & $(\mathit{bi}, \mathit{bj}, \mathit{best})$ \\
\midrule
1 & $5$ & $5$ & $5$ & $(1, 1, 5)$ \\
2 & $15$ & $20$ & $20$ & $(1, 2, 20)$ \\
3 & $-30$ & $-10$ & $0$，$\mathit{start} = 4$ & $(1, 2, 20)$ \\
4 & $10$ & $10$ & $10$ & $(1, 2, 20)$ \\
5 & $-5$ & $5$ & $5$ & $(1, 2, 20)$ \\
6 & $40$ & $45$ & $45$ & $(4, 6, 45)$ \\
7 & $10$ & $55$ & $55$ & $(4, 7, 55)$ \\
\bottomrule
\end{tabular}
\end{center}
回傳 $(4, 7, 55)$，也就是 $10, -5, 40, 10$，和 $55$，和題目給的答案相同。 全部是負數時，例如 $-3, -1$，
$\mathit{best}$ 一直是 $0$，回傳 $(1, 0, 0)$，即空的子序列。
\end{pastsoln}
